\section*{Chapter \ref{ch:g12:subduction-magmatism} --- Magmatism at Subduction Zones}

\begin{solution}{exo:g12:subduction-magmatism:1}
(a) a \omterm{def:g12:subduction-magmatism:arc}{volcanic arc}; (b) \omterm{def:g12:subduction-magmatism:hydrous-melting}{hydrous melting} (flux melting); (c) a \omterm{def:g12:subduction-magmatism:granodiorite}{granodiorite};
(d) a \omterm{def:g12:subduction-magmatism:batholith}{batholith}; (e) \omterm{def:g12:subduction-magmatism:continental-growth}{continental growth}.
\end{solution}

\begin{solution}{exo:g12:subduction-magmatism:2}
(1) At least $24/0.25 = \qty{96}{km/h}$, that is $\num{96000}/3600 \approx
\qty{27}{m/s}$. (2) The blast was at least $480/96 = 5$ times faster.
\end{solution}

\begin{solution}{exo:g12:subduction-magmatism:3}
(1) $838/1214 \approx 0.690$, about \qty{69}{\percent}. (2) On the side of
the plate that overrides, the one that does not sink. (3) Between the trench
and the front the \omterm{def:g11:subduction-zones:slab}{slab} is not yet deep enough: it has not reached the depth,
about a hundred kilometres, where its water is set free into a mantle hot
enough to melt once wet.
\end{solution}

\begin{solution}{exo:g12:subduction-magmatism:4}
(1) $45/(45 + 15) = 0.75$: \qty{75}{\percent} of the feldspar is
plagioclase. (2) Yes: it is coarse grained, holds more than
\qty{20}{\percent} quartz, and its plagioclase fraction lies between 65 and
\qty{90}{\percent}.
\end{solution}

\begin{solution}{exo:g12:subduction-magmatism:5}
Arc magmas are richer in silica, so they are more viscous; and they carry
the water that made them, so they are rich in gas. A high viscosity and a
high gas content both favour explosive eruptions.
\end{solution}

\begin{solution}{exo:g12:subduction-magmatism:6}
(1) $H = 150 \times \tan\ang{35} \approx 150 \times 0.700 \approx
\qty{105}{km}$. (2) It is the world average.
\end{solution}

\begin{solution}{exo:g12:subduction-magmatism:7}
$d = H/\tan\theta$: $105/\tan\ang{20} \approx 105/0.364 \approx
\qty{288}{km}$ at \ang{20}, and $105/\tan\ang{45} = \qty{105}{km}$ at
\ang{45}. The gentler the \omterm{def:g11:subduction-zones:slab}{slab}, the farther from the trench the arc stands.
\end{solution}

\begin{solution}{exo:g12:subduction-magmatism:8}
(1) $T = -5.104 \times 3.16^2 + 132.9 \times 3.16 + 1120.7 \approx -51.0 +
420.0 + 1120.7 \approx \qty{1490}{\celsius}$. (2) By about $1490 - 1000 =
\qty{490}{\celsius}$. (3) Dry, no: \qty{1300}{\celsius} is below
\qty{1490}{\celsius}. Wet, yes: it is above \qty{1000}{\celsius}.
\end{solution}

\begin{solution}{exo:g12:subduction-magmatism:9}
(1) $2.5/2.7 \approx 0.93$ of a year, about 340 days. (2) The avalanche moved
old rock that already existed; arcs make new rock from the mantle. Only the
volumes are compared.
\end{solution}

\begin{solution}{exo:g12:subduction-magmatism:10}
(1) $2.7 \times \num{1000000} = \qty{2.7e6}{km^3}$. (2) $\num{9e9}/2.7
\approx \num{3.3e9}$ years, 3.3 billion years. (3) It is less than the age of
the Earth, so arcs could have made the whole \omterm{def:g12:continental-crust:continental-crust}{continental crust}; but
\omterm{def:g11:subduction-zones:subduction}{subduction} also takes crust back, so in fact much of the crust has been
made and destroyed more than once.
\end{solution}

\begin{solution}{exo:g12:subduction-magmatism:11}
(1) Shallow earthquakes also occur in the overriding plate and on the fault
between the two plates; only the deep ones lie in the \omterm{def:g11:subduction-zones:slab}{slab} alone.
(2) $116 - 80 = \qty{36}{km}$ over \qty{86.4}{km}: $\tan\theta = 36/86.4
\approx 0.417$, $\theta \approx \ang{22.6}$.
\end{solution}

\begin{solution}{exo:g12:subduction-magmatism:12}
In the right triangle of the proposition, the \omterm{def:g11:subduction-zones:slab}{slab} is the hypotenuse $L$
and $d$ the side next to the angle $\theta$, so $\cos\theta = d/L$ and $L =
d/\cos\theta$. For Japan, $L = 270/\cos\ang{20} \approx 270/0.940 \approx
\qty{287}{km}$. At \qty{83}{mm} a year: $\num{2.87e8}/83 \approx \num{3.5e6}$
years, about 3.5 million years.
\end{solution}

\begin{solution}{exo:g12:subduction-magmatism:13}
(1) $V = \pi r^2 h = \pi \times 0.5^2 \times 8 \approx \qty{6.3}{km^3}$.
(2) About $6.3/2.5 \approx 2.5$ times the volume of the avalanche.
(3) $P = \rho g z$: at \qty{6}{km}, $2830 \times 9.81 \times 6000 \approx
\qty{1.67e8}{Pa} = \qty{0.17}{GPa}$; at \qty{14}{km}, about
\qty{3.89e8}{Pa} = \qty{0.39}{GPa}.
\end{solution}

\begin{solution}{exo:g12:subduction-magmatism:14}
(1) $2.7 - 3.2 = \qty{-0.5}{km^3}$ a year, a loss of \qty{5e5}{km^3} per
million years. (2) $500 \times \num{5e5} = \qty{2.5e8}{km^3}$, that is
$\num{2.5e8}/\num{9e9} \approx \qty{2.8}{\percent}$ of the crust. (3) Each
rate is an estimate, uncertain by much more than the difference of
\qty{0.5}{km^3} a year between them: within the errors, the two balance.
\end{solution}

\begin{solution}{exo:g12:subduction-magmatism:15}
(1) $\num{20000} \times 10 = \qty{2e5}{km^3}$. (2) $\num{2e5}/1 =
\num{200000}$ years for all the continental arcs of the world together. One
arc receives only a small part of that magma, so a single \omterm{def:g12:subduction-magmatism:batholith}{batholith} takes
far longer: it is built \omterm{def:g11:igneous-rocks:intrusion}{pluton} by \omterm{def:g11:igneous-rocks:intrusion}{pluton}, over millions of years. (3) At
least the \qty{5}{km} of rock above its top.
\end{solution}

\begin{solution}{pb:g12:subduction-magmatism:1}
\textbf{1.} $111.2 \times \cos\ang{39} \approx 111.2 \times 0.777 \approx
\qty{86.4}{km}$.

\textbf{2.} $(144.1 - 141.0) \times 86.4 \approx 3.1 \times 86.4 \approx
\qty{268}{km}$, about \qty{270}{km}.

\textbf{3.} $(141.0 - 139.7) \times 86.4 \approx \qty{112}{km}$.

\textbf{4.} On the overriding plate, the one that does not sink, west of the
trench: the Pacific Plate sinks beneath it.

\textbf{5.} There, the \omterm{def:g11:subduction-zones:slab}{slab} is not yet deep enough to release its water into
a mantle hot enough to melt.

\textbf{6.} $116 - 80 = \qty{36}{km}$.

\textbf{7.} One degree of longitude, \qty{86.4}{km}.

\textbf{8.} $\tan\theta = 36/86.4 \approx 0.417$, so $\theta \approx
\ang{22.6}$.

\textbf{9.} The line rises \qty{36}{km} every \qty{86.4}{km} eastwards, so it
reaches the surface $98/0.417 \approx \qty{235}{km}$ east of Iwate, at
$141.0 + 235/86.4 \approx \ang{143.7}$~E: about $(144.1 - 143.7) \times 86.4
\approx \qty{33}{km}$ west of the trench.

\textbf{10.} The \omterm{def:g11:subduction-zones:slab}{slab} is not a flat plane: it enters the trench almost
horizontal and bends down more and more as it sinks.

\textbf{11.} $268 \times \tan\ang{20} \approx 268 \times 0.364 \approx
\qty{98}{km}$.

\textbf{12.} They agree: \qty{98}{km} both ways.

\textbf{13.} Both lie within the world range, a little shallower than the
mean of \qty{105}{km}.

\textbf{14.} Megata lies \qty{112}{km} west of Iwate, so the \omterm{def:g11:subduction-zones:slab}{slab} beneath it
is about $98 + 0.417 \times 112 \approx \qty{145}{km}$ deep.

\textbf{15.} In the mantle, about $98 - 35 = \qty{63}{km}$ below the base of
the crust.

\textbf{16.} It entered near the ridge where the crust formed, from sea
water; much of it is held in serpentine.

\textbf{17.} $T = -5.104 \times 3.16^2 + 132.9 \times 3.16 + 1120.7 \approx
\qty{1490}{\celsius}$.

\textbf{18.} By $1490 - 1000 = 490$ to $1490 - 950 = \qty{540}{\celsius}$:
about \qty{500}{\celsius}.

\textbf{19.} The \omterm{def:g11:subduction-zones:slab}{slab} is several hundred degrees colder than the mantle
around it, too cold to melt even with water: it is the hot mantle wedge above
it, once wetted, that melts.

\textbf{20.} The \omterm{def:g11:subduction-zones:slab}{slab} lies about \textbf{\qty{100}{km} beneath the volcanic
front} of north-east Japan (\qty{98}{km} by both methods).
\end{solution}
